\section{Resolved 3D Comparison at \texorpdfstring{\(T=1\)}{T=1}}
\label{sec:resolved-3d-comparison}

This section records the numerical comparison produced by the Julia
implementation of the stenosis-aware 1D area--flow model against the available
resolved 3D velocity data. The comparison is an output-operator check, not a
clinical validation claim and not an exact cross-section quadrature of the 3D
tetrahedral solution. The 3D diagnostic averages node-centered axial velocity
over thin axial slabs and compares those values with the interpolated 1D
quantity \(Q/A\). Radial profiles use the same slab nodes binned by
\(\sqrt{x^2+y^2}/R_0(z)\) and compare them with the 1D closure profile.
The tables and figures in this section are baseline default Newtonian records;
non-Newtonian rheology descriptors require separately generated comparison
records.
\todo[inline]{ADVISOR[BLOCKER]: This numerical section needs a reproducibility
paragraph naming the exact computed realization: solver closure, grid \(N\),
timestep/CFL policy, inlet/outlet data, 3D source files or case IDs,
output-operator settings, and output-record location. Without that,
Tables~\ref{tab:t1-3d-comparison-summary}--\ref{tab:t1-radial-profile-summary}
are not auditable.}
\red{All numbers below come from the recorded Newtonian realization with
[insert grid, timestep policy, boundary data, source-data identifiers, and
output-record path], interpreted under
Convention~\ref{conv:computed-realization-record}.}
\todo[inline]{ADVISOR[BLOCKER]: The main model defines \(A\) as physical
cross-sectional area and \(Q=\int_S u_z\,dS\), while the solver appendix
switches to \(a=R^2\) with \(A_{\mathrm{phys}}=\pi a\). State whether the
solver's \(Q\) is physical volumetric flow or a \(\pi\)-scaled flow, and define
exactly whether comparison velocities use \(Q/A_{\mathrm{phys}}\), \(Q/a\), or
another recorded output.}
\todo{ADVISOR[MAJOR]: Declare the active 1D closure used for Section~2. The text
says variable-radius corrections are inactive unless a closure identifier is
declared, but the numerical-method appendix describes the current Canic
extended flux; the comparison reader needs to know which model generated
Tables~\ref{tab:t1-3d-comparison-summary}--\ref{tab:t1-radial-profile-summary}.}
\red{For the comparison below, the active closure is [insert closure
identifier]; the Canic variable-radius terms are [active or inactive] in the
recorded flux and source.}

The two included upstream states are case 77, treated as a 23\% stenosis, and
case 60, treated as a 40\% stenosis. Both XDMF files report
\(t=0.9994999999999453\), within \(5.0\times 10^{-4}\) of \(T=1\). Upstream
case 50 was not used because its XDMF time is
\(t=1.4994999999998904\), outside the \(10^{-3}\) tolerance for the \(T=1\)
comparison.

\begin{table}[!htb]
    \centering
    \scriptsize
    \caption{Section-mean and radial-profile discrepancy summary for the
    \(T=1\) resolved-node comparison. Velocity errors are in cm/s.}
    \label{tab:t1-3d-comparison-summary}
    \begin{tabular}{@{}rrrrrrrrrr@{}}
        \toprule
        \textbf{Case}
            & \textbf{Severity}
            & \textbf{Sections}
            & \textbf{Min nodes}
            & \textbf{Mean \(|e_s|\)}
            & \textbf{Max \(|e_s|\)}
            & \textbf{Mean rel.}
            & \textbf{Max rel.}
            & \textbf{Mean \(|e_r|\)}
            & \textbf{Max \(|e_r|\)} \\
        \midrule
        77 & 23\% & 200 & 64 & 7.471 & 26.26 & 0.4534 & 1.622 & 10.03 & 23.01 \\
        60 & 40\% & 200 & 43 & 9.440 & 47.48 & 0.5141 & 2.458 & 15.35 & 41.72 \\
        \bottomrule
    \end{tabular}
\end{table}

Figure~\ref{fig:t1-section-mean-comparison} shows the section-mean comparison.
The 1D model captures the smooth upstream and downstream trend but overpredicts
the highest-velocity region near the stenosis, especially for the 40\% case.
The largest section errors occur for case 60 around \(z=2.4\)--\(2.8\) cm,
where the 1D mean velocity is between 58 and 67 cm/s while the node-slab 3D
mean is between 19 and 27 cm/s.
\todo{ADVISOR[MAJOR]: Do not leave the reader to infer the meaning of relative
errors above one. State whether these errors are acceptable for the report's
purpose, expected from closure mismatch, or evidence that the current comparison
is diagnostic rather than validating.}
\red{Because the largest relative errors exceed unity near the stenosis, the
result is best read as a diagnostic of where the reduced closure and node-slab
output operator diverge, not as accuracy validation.}

\begin{figure}[!htb]
    \centering
    \captionsetup{aboveskip=2pt,belowskip=2pt}
    \figtikz{section-mean-comparison}
    \caption{Section-mean axial velocity comparison at \(T=1\). Solid curves
    are resolved 3D node-slab means; dashed curves are interpolated 1D
    \(Q/A\).}
    \label{fig:t1-section-mean-comparison}
\end{figure}

\begin{table}[!htb]
    \centering
    \scriptsize
    \caption{Five largest section-mean absolute errors across the two
    included \(T=1\) cases. Velocities and errors are in cm/s.}
    \label{tab:t1-largest-section-errors}
    \begin{tabular}{@{}rrrrrrrr@{}}
        \toprule
        \textbf{Rank}
            & \textbf{Case}
            & \(\boldsymbol{z}\) \textbf{(cm)}
            & \(\boldsymbol{u_{\mathrm{1D}}}\)
            & \(\boldsymbol{u_{\mathrm{3D}}}\)
            & \(\boldsymbol{|e|}\)
            & \textbf{Rel.}
            & \textbf{Nodes} \\
        \midrule
        1 & 60 & 2.714 & 66.79 & 19.31 & 47.48 & 2.458 & 72 \\
        2 & 60 & 2.774 & 63.74 & 24.57 & 39.17 & 1.594 & 60 \\
        3 & 60 & 2.412 & 61.90 & 23.51 & 38.39 & 1.633 & 67 \\
        4 & 60 & 2.563 & 63.71 & 27.15 & 36.56 & 1.347 & 72 \\
        5 & 60 & 2.834 & 57.59 & 21.68 & 35.92 & 1.657 & 82 \\
        \bottomrule
    \end{tabular}
\end{table}

The radial-profile diagnostics in
Figure~\ref{fig:t1-radial-profile-comparison} show the same pattern at the
profile level. For both cases, the throat slice \(z=2.451\) cm has fewer
populated radial bins because fewer 3D nodes lie in the narrowest cross-section
slab. The largest radial-profile discrepancy is downstream of the throat at
\(z=2.951\) cm, where case 60 reaches a maximum per-bin absolute error of
41.72 cm/s.
\todo[inline]{ADVISOR[MAJOR]: Define the radial closure profile used for
Figure~\ref{fig:t1-radial-profile-comparison}, including its formula,
normalization to the local mean velocity, and relation to \(\alpha=1.1\).
Without that definition, the radial-profile discrepancy is not reproducible
from the manuscript's \(A,Q,\bar u\) variables.}

\begin{figure}[!htb]
    \centering
    \captionsetup{aboveskip=2pt,belowskip=2pt}
    \figtikz{radial-profile-comparison}
    \caption{Radial velocity profile comparison at the default throat-centered
    slices. Markers are resolved 3D radial-bin means; dashed curves are the 1D
    closure profile reconstructed from the local \(Q/A\).}
    \label{fig:t1-radial-profile-comparison}
\end{figure}

\begin{table}[!htb]
    \centering
    \scriptsize
    \caption{Radial-profile discrepancy by case and slice. Errors are in
    cm/s.}
    \label{tab:t1-radial-profile-summary}
    \begin{tabular}{@{}rrrrr@{}}
        \toprule
        \textbf{Case}
            & \(\boldsymbol{z}\) \textbf{(cm)}
            & \textbf{Populated bins}
            & \textbf{Mean \(|e_r|\)}
            & \textbf{Max \(|e_r|\)} \\
        \midrule
        77 & 1.951 & 17 & 9.404 & 18.60 \\
        77 & 2.451 & 12 & 6.949 & 13.33 \\
        77 & 2.951 & 16 & 13.01 & 23.01 \\
        60 & 1.951 & 18 & 10.04 & 18.90 \\
        60 & 2.451 & 10 & 7.562 & 17.17 \\
        60 & 2.951 & 17 & 25.55 & 41.72 \\
        \bottomrule
    \end{tabular}
\end{table}

These discrepancies should be read in light of the comparison boundary in
Section~\ref{subsubsec:model-hierarchy}. The present 1D solver is compared to
a resolved 3D velocity field through simple node-slab output operators. The
result is useful for locating where the reduced velocity closure departs from
the resolved state, but it does not by itself identify whether the discrepancy
comes from the 1D closure, the boundary data, the 3D post-processing operator,
or differences between the numerical formulations.
\todo[inline]{ADVISOR[BLOCKER]: The report ends the main text with a limitation
caveat and then enters the appendices. Add a conclusion that states what the
\(T=1\) comparison shows, what it does not show, and the next defensible study
step.}
\red{In summary, the \(T=1\) comparison supports using the 1D solution as a
controlled reduced-output diagnostic, but the large stenosis-region velocity
errors mean the present record should be read as a localization of
closure/output-operator mismatch rather than validation of the 1D model.}
